\begin{table}[H]
\centering\footnotesize
\setlength{\tabcolsep}{4pt}
\caption{FPM recebido e população: Barreiras e Camaçari (2014--2025)}
\label{tab:fpm}
\begin{adjustbox}{max width=\textwidth}
\begin{tabular}{lrrrrr}
\toprule
Ano & FPM Barreiras & FPM Camaçari & Barreiras/Camaçari & Pop. Barreiras & Pop. Camaçari \\
\midrule
2014 & 73,9 & 75,7 & 97,5\% & 152.208 & 281.413 \\
2015 & 80,3 & 82,9 & 96,9\% & 153.918 & 286.919 \\
2016 & 93,6 & 96,1 & 97,4\% & 155.519 & 292.074 \\
2017 & 90,3 & 93,2 & 96,9\% & 157.638 & 296.893 \\
2018 & 95,6 & 95,6 & 100,0\% & 153.831 & 293.723 \\
2019 & 103,5 & 103,5 & 100,0\% & 155.439 & 299.132 \\
2020 & 98,9 & 98,9 & 100,0\% & 156.975 & 304.302 \\
2021 & 132,1 & 132,1 & 100,0\% & 158.432 & 309.208 \\
2022 & 166,6 & 166,6 & 100,0\% & 159.734 & 300.372 \\
2023 & 170,7 & 170,7 & 100,0\% & 165.200 & 309.883 \\
2024 & 196,1 & 196,1 & 100,0\% & 170.667 & 319.394 \\
2025 & 214,1 & 214,1 & 100,0\% & 171.634 & 321.636 \\
\bottomrule
\end{tabular}
\end{adjustbox}
\par\vspace{3pt}\parbox{\textwidth}{\scriptsize Nota: FPM em R\$ milhões correntes, todas as parcelas (Siconfi, DCA). População: estimativas anuais do IBGE, Censo 2022 em 2022 e interpolação em 2023. O coeficiente do FPM de cada ano é fixado pelo TCU com a estimativa populacional do ano anterior; a Lei Complementar n.º 165/2019 manteve, até o censo seguinte, os coeficientes de 2018 dos municípios que perderiam coeficiente por causa das estimativas anuais. Fonte: elaboração própria com dados do Siconfi e do IBGE.}
\end{table}
